% Part: first-order-logic
% Chapter: introduction
% Section: first-order-logic

\documentclass[../../../include/open-logic-section]{subfiles}

\begin{document}

\olfileid{fol}{int}{fol}

\olsection{प्रथम-क्रम तर्कशास्त्र}

आकारिक तर्कशास्त्राशी तुमची पहिली ओळख झाली तेव्हापासून प्रथम-क्रम
तर्कशास्त्र बहुधा तुम्हाला परिचित असेल.\footnote{खरे तर, तुम्ही त्याच्याशी
परिचित आहात असे आम्ही जवळजवळ गृहीतच धरतो!{} तसे नसल्यास,
\emph{forall x} \citep{Magnus2021} यासारख्या अधिक प्राथमिक
पाठ्यपुस्तकाची उजळणी करू शकता.} ते तुम्हाला ``संख्यापक तर्कशास्त्र''
किंवा ``विधेय तर्कशास्त्र'' या नावांनी माहीत असू शकते. सर्वप्रथम,
प्रथम-क्रम तर्कशास्त्र ही एक आकारिक भाषा आहे. म्हणजे तिचा एक विशिष्ट
शब्दसंग्रह असतो आणि तिच्या अभिव्यक्ती या त्या शब्दसंग्रहातील चिन्हमाला
असतात. पण प्रत्येक चिन्हमाला अनुमत नसते. अनुमत अभिव्यक्तींचे विविध
प्रकार आहेत: पदे, !!{formula}s आणि !!{sentence}s. आपल्याला मुख्यतः
प्रथम-क्रम तर्कशास्त्रातील !!{sentence}s मध्ये रस आहे: ती आपल्याला
इंग्रजीतील वाक्यांचे आकारिक प्रतिरूप देतात आणि त्यांच्याविषयी
तर्कशास्त्रज्ञाला नेहमी रस असणारे प्रश्न आपण विचारू शकतो. उदाहरणार्थ:
\begin{itemize}
    \item $!B$ हे~$!A$ पासून तार्किकरीत्या निष्पन्न होते का?
    \item $!A$ तार्किकरीत्या सत्य, तार्किकरीत्या असत्य, की
परिस्थितीवश आहे?
    \item $!A$ आणि $!B$ सममूल्य आहेत का?
\end{itemize}

हे प्रश्न मुख्यतः प्रथम-क्रम तर्कशास्त्रातील !!{sentence}s च्या
``अर्था''विषयीचे आहेत. उदाहरणार्थ, $!B$ हे~$!A$ पासून तार्किकरीत्या
निष्पन्न होते का या प्रश्नाचे तत्त्वज्ञ पुढीलप्रमाणे विश्लेषण करेल:
$!A$ सत्य पण~$!B$ असत्य असणारा काही प्रसंग आहे का ($!B$ हे~$!A$
पासून निष्पन्न होत नाही), की $!A$ सत्य करणारा प्रत्येक प्रसंग~$!B$ ही
सत्य करतो ($!B$ हे~$!A$ पासून निष्पन्न होते)? पण ``प्रसंग'' म्हणजे काय
हे अजून सांगितलेले नाही---ते \emph{चिन्हार्थमीमांसेचे} काम आहे.
प्रथम-क्रम तर्कशास्त्राची चिन्हार्थमीमांसा तत्त्वज्ञाच्या ``प्रसंग'' या
अंतःप्रज्ञ कल्पनेचे गणितीयदृष्ट्या काटेकोर प्रतिमान देते आणि---हे
महत्त्वाचे आहे---एखाद्या प्रसंगात !!a{sentence}~$!A$ \emph{सत्य असणे}
म्हणजे काय याचेही प्रतिमान देते. आपण विकसित करणार असलेल्या या
गणितीयदृष्ट्या काटेकोर प्रतिमानाला !!a{structure} म्हणतो. ``यात सत्य''
याला काटेकोर अर्थ देणाऱ्या संबंधाला \emph{पूर्ति}-संबंध म्हणतात. म्हणून
!!{sentence}s~$!A$ आणि !!{structure}s~$\Struct{M}$ यांच्यासाठी आपण
``$!A$ ची~$\Struct{M}$ मध्ये पूर्ति होते'' (चिन्हांत:
$\Sat{M}{!A}$) हे परिभाषित करणार आहोत. हे झाल्यावर ``पासून निष्पन्न
होते'' किंवा ``तार्किकरीत्या सत्य आहे'' यांसारख्या इतर चिन्हार्थविषयक
संज्ञांच्याही काटेकोर व्याख्या देता येतील. उदाहरणार्थ,
$\lforall[x][(!A(x) \lif !B(x))],
\lexists[x][!A(x)] \Entails \lexists[x][!B(x)]$ आहे की नाही हे पुन्हा
गणितीय काटेकोरपणे ठरवणे या व्याख्यांमुळे शक्य होईल. अर्थातच उत्तर
``होय'' असेल. इंग्रजी वाक्यांचे प्रथम-क्रम तर्कशास्त्रात चिन्हांकन
करण्याचे प्रशिक्षण तुम्ही आधीच घेतले असेल, तर ही, उदाहरणार्थ, ``सर्व
मुंग्या कीटक आहेत, मुंग्या अस्तित्वात आहेत, म्हणून कीटक अस्तित्वात
आहेत'' यांची चिन्हांकने आहेत हे तुम्ही ओळखाल. हा उघडपणे वैध युक्तिवाद
आहे आणि म्हणून आपल्या आकारिक भाषेतील ``पासून निष्पन्न होते'' याच्या
गणितीय प्रतिमानानेही तेच उत्तर दिले पाहिजे.
%READERNOTE{OLFOL-001: गोठवलेल्या इंग्रजी मजकुरातील या तार्किक निष्पन्नता-सूत्रात आणि पुढील दोन पुनरुक्तींमध्ये सार्वत्रिक सूत्राचा बंद करणारा चौकोनी कंस चुकीच्या जागी आहे किंवा गहाळ आहे, तर निष्कर्षानंतर एक जादा कंस आहे. मराठी मजकुरात तिन्ही ठिकाणी समान, सुव्यवस्थित आधारविधाने आणि निष्कर्ष दिले आहेत.}

आकारिक तर्कशास्त्राशी तुमची पहिली ओळख झाली तेव्हापासून आठवत असणारा
आणखी एक विषय म्हणजे \emph{!!{derivation}s}. तुम्ही आकारिक
तर्कशास्त्राचा पहिला अभ्यासक्रम घेतला असेल, तर तुमच्या शिक्षकांनी अशा
!!{derivation}s शोधण्याचा सराव तुमच्याकडून करून घेतला असेल; कदाचित
वरील तार्किक निष्पन्नता लागू असल्याचे दाखवणारी !!a{derivation} देखील.
!!{derivation}s देण्याच्या अनेक पद्धती आहेत: तुम्ही ``नैसर्गिक निगमन''
किंवा ``सत्यता-वृक्ष'' नावाची पद्धत केली असेल; पण इतर अनेक पद्धतीही
आहेत. वरील तर्कशास्त्रीय प्रश्नांची उत्तरे देता येतील अशी साधने पुरवणे
हे !!{derivation} पद्धतींचे उद्दिष्ट आहे. उदाहरणार्थ, ज्या नैसर्गिक
निगमनातील !!{derivation} मध्ये $\lforall[x][(!A(x) \lif !B(x))]$ आणि
$\lexists[x][!A(x)]$ ही आधारविधाने आणि $\lexists[x][!B(x)]$ हा
निष्कर्ष (शेवटची ओळ) आहे, ती $\lexists[x][!B(x)]$ हे
$\lforall[x][(!A(x) \lif !B(x))]$ आणि $\lexists[x][!A(x)]$ यांपासून
तार्किकरीत्या निष्पन्न होते याची \emph{पुष्टी करते}.

पण असे का? वरवर पाहता, !!{derivation} पद्धतींचा चिन्हार्थमीमांसेशी
काहीच संबंध नाही: आकारिक !!{derivation} देणे म्हणजे केवळ काही
नियमांनी नियंत्रित रीतीने चिन्हांची मांडणी करणे; त्यात ``प्रसंगां''चा
किंवा ``यात सत्य''चा मुळीच उल्लेख नसतो. !!{derivation} पद्धती आणि
चिन्हार्थमीमांसा यांच्यातील संबंध अधितार्किक अन्वेषणाने प्रस्थापित
करावा लागतो. उदाहरणार्थ, आधारविधाने
$\lforall[x][(!A(x) \lif !B(x))]$ आणि $\lexists[x][!A(x)]$ यांपासून
$\lexists[x][!B(x)]$ ची आकारिक !!{derivation} शक्य असते तेव्हा आणि
केवळ तेव्हाच $\lforall[x][(!A(x) \lif !B(x))]$ आणि
$\lexists[x][!A(x)]$ ही दोन्ही मिळून
$\lexists[x][!B(x)]$ तार्किकरीत्या निष्पन्न करतात, याची गणितीय
सिद्धता आवश्यक आहे. मात्र, हे करता येण्यापूर्वी विन्यासमीमांसा आणि
चिन्हार्थमीमांसा यांच्या व्याख्या अचूक करण्याचे बरेच कष्टप्रद काम
करावे लागते.

\end{document}
